\documentclass{article}
\usepackage{amsmath, amsthm, amssymb, amsfonts, mathrsfs,relsize, thmtools, mathtools, stmaryrd, mlmodern, graphicx, setspace, geometry, float, quiver, xspace, hyperref, framed, tcolorbox, mathrsfs, enumerate}
\usepackage[T1]{fontenc}
\usepackage[all]{xy}
\usepackage[utf8]{inputenc}
\usepackage[english]{babel}
\usepackage[dvipsnames]{xcolor}
\usepackage[document]{ragged2e}

\xyoption{curve}
\xyoption{2cell}
\UseAllTwocells

% other symbols
\newcommand{\isom}{\cong}
\DeclareMathOperator{\idem}{E}
\DeclareMathOperator{\supp}{supp}
\DeclareMathOperator{\End}{End}
\DeclareMathOperator{\Inc}{Inc}
\DeclareMathOperator{\img}{im}
\DeclareMathOperator{\jac}{J}
\DeclareMathOperator{\fib}{\textit{P} \colon \mathcal{E} \longrightarrow \mathcal{B}}
\DeclareMathOperator{\opfib}{\textit{P} \colon \mathcal{E}^{op} \longrightarrow \mathcal{B}^{op}}
\DeclareMathOperator{\comp}{\circ}
\DeclareMathOperator{\mor}{\longrightarrow}
\DeclareMathOperator{\is}{\equiv}
\DeclareMathOperator{\forw}{\Longrightarrow}
\DeclareMathOperator{\back}{\Longleftarrow}
\DeclareMathOperator{\adj}{\dashv}
\DeclareMathOperator{\2cell}{\Longrightarrow}



\newcommand{\dfib}[3]{\cat{#1} \colon \mathbb{#2} \longrightarrow \mathbb{#3}}
\newcommand{\compdfib}[4]{\cat{#1}_{#4} \colon \mathbb{#2}_{#4} \longrightarrow \mathbb{#3}_{#4}}
\newcommand*{\dcat}[1]{\mathbb{#1}}
\newcommand*{\bspan}[1]{\mathbb{S}\text{pan}(#1)}
\newcommand*{\slice}[1]{\mathcal{C}/#1}
\newcommand*{\slicec}[1]{\mathbf{Fib}/#1}
\newcommand*{\dspan}[1]{\mathbb{S}\mathbf{pan}(#1)}
\newcommand{\cat}[1]{\mathcal{#1}}
\newcommand{\hspan}[5]{\xymatrix@=.7cm{#1 & #2 \ar[l]_{#4}  \ar[r]^{#5} & #3}}
\newcommand{\cospan}[5]{\xymatrix@=.7cm{#1 \ar[r]^{#4} & #2 & #3 \ar[l]_{#5} & }}
\newcommand{\proar}[2]{\xymatrix@=.7cm{#1 \ar[r]|-{\shortmid} & #2}}
\newcommand{\inproar}[3]{\xymatrix@=.7cm{#2 \ar[r]|-{\shortmid}^{#1} & #3}}
\newcommand{\lproar}[2]{\xymatrix@=.7cm{#1 \ar[l]|-{\shortmid} & #2}}
\newcommand{\pushoutcorner}[1][dr]{\save*!/#1+1.2pc/#1:(1,-1)@^{|-}\restore}
\newcommand{\pullbackcorner}[1][dr]{\save*!/#1-1.2pc/#1:(-1,1)@^{|-}\restore}
\newcommand{\chapquote}[2]{\begin{quotation} \textit{#1} \end{quotation} \begin{flushright} - #2 \end{flushright} }
\newcommand{\cleav}[2]{\varphi(#1,#2)}
\newcommand{\ter}[1]{1_{\mathcal{#1}}}

\def\ie{i.e.\@\xspace}




\colorlet{Gray}{White!90!Periwinkle}
\colorlet{Orange}{Orange!40}
\colorlet{Green}{Green!40}
\colorlet{Blue}{Blue!40}
\colorlet{Goldenrod}{Goldenrod!40}

\newcommand{\HRule}[1]{\rule{\linewidth}{#1}}

\declaretheoremstyle[name=Theorem,]{thmsty}
\declaretheorem[style=thmsty,numberwithin=section]{theorem}
\tcolorboxenvironment{theorem}{colback=Gray}

\declaretheoremstyle[name=Proposition,]{prosty}
\declaretheorem[style=prosty,numberlike=theorem]{proposition}
\tcolorboxenvironment{proposition}{colback=Orange}

\declaretheoremstyle[name=Definition,]{definition}
\declaretheorem[style=definition,numberwithin=section]{definition}
\tcolorboxenvironment{definition}{colback=Green}

\declaretheoremstyle[name=Remark,]{remark}
\declaretheorem[style=remark,numberlike=theorem]{remark}
\tcolorboxenvironment{remark}{colback=Blue}

\declaretheoremstyle[name=Example,]{example}
\declaretheorem[style=example,numberlike=theorem]{example}
\tcolorboxenvironment{example}{colback=Cyan}

\declaretheoremstyle[name=Lemma,]{lemma}
\declaretheorem[style=lemma,numberlike=theorem]{lemma}
\tcolorboxenvironment{lemma}{colback=Goldenrod}



\def\ie{i.e.\@\xspace}



\setstretch{1.2}
\geometry{
    textheight=9in,
    textwidth=5.5in,
    top=1in,
    headheight=12pt,
    headsep=25pt,
    footskip=30pt
}

% ------------------------------------------------------------------------------

\begin{document}

% ------------------------------------------------------------------------------
% Cover Page and ToC
% ------------------------------------------------------------------------------

\title{ \normalsize \textsc{}
		\\ [2.0cm]
		\HRule{1.5pt} \\
		\LARGE \textbf{\uppercase{(Co)Completeness of Fibred Double Categories}
		\HRule{2.0pt} \\ [0.6cm] \LARGE{PhD Writings} \vspace*{10\baselineskip}}
		}
\date{}
\author{\textbf{Author} \\ 
		Siyabonga Mthimkhulu \\
		Muizenberg, Cape Town}

\maketitle
\newpage

\tableofcontents
\newpage

% ------------------------------------------------------------------------------
These are edited notes from \cite{TS} including proofs that are originally not found in the text. \\ They cover accordingly Internal Sums and Products.


\section{Depend Sums and Products}

Given a category $\mathcal{C}$ and a morphism morphism $f \colon X \mor Y$ in $\mathcal{C}$ s.t pullbacks along $f$ exist. We say

\begin{definition}
	A depend sum along $f$ is the functor $\sum_{f} \colon            \mathcal{C}/X \mor \mathcal{C}/Y$ that is left adjoint to the change-of-base functor $f^\ast \colon \mathcal{C}/Y \mor \mathcal{C}/X$ 
	\[\xymatrix@=0.5cm{
		\slice{X} \ar@/^1.5pc/[rrrr]^{\sum_f} &&\perp  && \slice{Y} \ar@/^1.5pc/[llll]^{f^\ast} 
	}\]
	
\end{definition}

Dually 

\begin{definition}
	A dependent product along $g \colon X \mor Y$ is, if it exists, the right adjoint functor $\prod_{g} \colon \slice{X} \mor \slice{Y}$ to the change-of-base along $g$. 
	
	\[\xymatrix@=0.5cm{
		\slice{X} \ar@/^1.5pc/[rrrr]^{\prod_f} &&\top  && \slice{Y} \ar@/^1.5pc/[llll]^{g_\ast} 
	}\]
\end{definition}

\begin{remark}
	A category with dependent products necessarily has pullbacks. Thus we have 
	
	\begin{equation*}
		\sum_f \adj f^\ast \adj \prod_f
	\end{equation*}
	
\end{remark}


\[\xymatrix@=1.5cm{
	X \rtwocell^{f}_{g}{\eta} & B
}
\]


\section{Opfibrations}
Let $\fib$ be a functor. 

\begin{enumerate}
	\item  A morphism $\varphi \colon X \mor Y$ in $\mathcal{E}$ over $u \colon I \mor J$ is opcartesian if it is cartesian over $u$ for $P \colon \mathcal{E}^{op} \mor \mathcal{B}^{op}$. That is, if each morphism 
	$\psi \colon X \mor Z$ with $P\psi = v \comp u$ in $\mathcal{B}$, uniquely determines a morphism $\theta \colon Y \mor Z$ over $v$. \\
	
	\item A functor $\fib$ is an opfibration if $\opfib$ is a fibration. Equivalently, if for each $u \colon I \mor J$ in $\mathcal{B}$, there is an opcartesian morphism $\varphi \colon X \mor Y$ over $u$. \\
	
	\item  A bifibration is a functor $\fib$ that is a fibration and an opfibration.
\end{enumerate}

\begin{remark}
	Usually it is written $u^*X \mor X $ for a cartesian lifting of $u \colon I \mor J$. Similarly, it also written $X \mor u_{!}X$ or $X \mor \coprod_uX$ for opcartesian lifting of $u$. \\
	
	The assignment $X \mapsto u_{!}X$ extends to a functor which may be called an extension --- because it turns out to be left adjoint to $u$.
\end{remark} 


\begin{example}
	The domain functor $\partial_0 \equiv \mbox{dom} \colon \mathcal{B}^2 \mor \mathcal{B}$ is an opfibration if and only if $\mathcal{B}$ has pushouts. Pushouts are precisely the opcartesian morphisms in $\mathcal{B}^2$. \\
	The fibre over $I \in \mathcal{B}$ is the opslice $I/\mathcal{B}$. 
\end{example}

\begin{lemma}[B. Jacobs]
	A fibration is a bifibratoin (\ie additionally it is an opfibration) if and only if each reindexing functor $u^*$ has a left adjoint, denoted $\coprod_u$. 
\end{lemma}

\begin{proof}
	Let $\fib$ be a fibration. For a morphism $u \colon I \mor J$ in $\mathcal{B}$ and objects $X \in \mathcal{E}_I $ and $Y \in \mathcal{E}_J$ consider the following chain
	\[
	\mathcal{E}_J\left(\coprod_uX,Y\right) \overset{*}{\isom} \mathcal{E}_I(X, u^*Y) \isom \mathcal{E}_u(X,Y) \overset{**}{\isom} \mathcal{E}_J\left(\coprod_uX,Y\right)
	\]
	where $\isom$ is due to $P$ being a fibration. Then 
	\begin{align*}
		\mbox{left adjoints $\coprod_u$ exist} &\iff  \overset{*}{\isom} \mbox{ exist}\\
												&\iff \overset{**}{\isom} \mbox{ exist} \\
												&\iff P \mbox{ is an opfibrati
													on.}
	\end{align*}
\end{proof}

\begin{remark}
	The Beck-Chevalley condition is not required for these $\coprod_u$'s in the above lemma. However, this result implies that every fibration with sums (\ie adjunctions $\coprod_u \adj u^*$ satisying Beck-Chevalley) is a bifibration. Thus, codomain fibrations are bifibrations.
\end{remark}

\section{Geometric Morphisms as Fibrations}

\begin{theorem}
	Given a functor $F \colon \mathcal{C} \mor \mathcal{D}$ with a fully faithful right adjoint $G$ s.t $F$ preserves finite limits and $\mathcal{C}$ is finitely complete, then $F$ is a Street fibration.
\end{theorem}

Before we prove the above theorem, we provide this lemma
\begin{lemma}
	If a functor $G \colon \mathcal{D} \mor \mathcal{C}$ has a left adjoint $F$ with counit $\varepsilon$. Then $G$ is fully faithful iff the $\varepsilon$ is an isomorphism.
\end{lemma}  

\begin{proof}\label{counit}
	Consider
	\[\xymatrix@=0.5cm{
		\mbox{Hom}(X,X') \ar[rr]^{G} \ar[ddrr]_{\varepsilon_X}  && \mbox{Hom}(GX,GX')	\ar[dd]_{\thicksim}^{\mbox{ad}} \\
		\\
		&&\mbox{Hom}(FGX,X')
	}
	\]
	then $\mbox{Hom}(X,X') \xlongrightarrow{\thicksim} \mbox{Hom}(GX,GX')$ iff $\varepsilon_X$ is a bijection. $G$ is then fully faithful iff $FGX \mor X$ is an isomorphism for all $X$. 
\end{proof}

We can now prove the theorem
\begin{proof}
	Given $Y$ in $\mathcal{C}$ with $F(Y) = J$ and a morphism $u \colon I \to J$, let $u^*Y$ be the pullback of the unit $\eta \colon Y \mor G(J) = GF(Y)$ along $G(u) \colon G(I) \mor G(J)$. Since $F(\eta_Z)$ is an isomorphism (as $G$ is fully faithful) and $F$ preserves pullbacks, $F$ also sends the morphism $u^*(Y) \mor G(I)$ to an isomorphism. As illustrated in the diagram below
	
	\[\xymatrix@=0.5cm{
		u^*Y \ar[rr]^{\varphi} \ar[dd] && Y \ar[dd]^{\eta_Y} \\
		\\
		G(I) \ar[rr]_{G(u)} && G(J)
	}
	\]
	Claim: $\varphi \colon u^*Y \mor Y$ is cartesian. \\
	
	Suppose $\psi \colon Z \mor Y$ is a morphism over $v \colon K \mor I$ and $\psi$ is over $u \comp v$. Then we have $Gu \comp Gv \comp \eta_Z = \eta_Y \comp \psi$ by the naturality of the adjunction unit. The universal property of the pullback gives a  unique morphism $\theta \colon Z \mor u^*Y$ over $v$. As shown below
	
	\[\xymatrix@=0.5cm{
		Z \ar@{-->}[rr]^{\theta} \ar[dd]_{\eta_Z} \ar@/^1.5pc/[rrrr]^{\psi} && u^*Y \ar[rr]^{\phi} \ar[dd] && Y \ar[dd]^{\eta_Y} \\
		\\
		G(K) \ar[rr]_{G(v)} && G(I) \ar[rr]_{G(u)} && G(J) \\
		\\
		K \ar[rr]^{v} && I \ar[rr]^{u} && J
	}
	\]
	Furthermore, for any other morphism $\tilde{\theta} \colon Z \mor u^*Y$ over $v$ that commutes with $\psi$ and $\varphi$, it must satisfy $G(v) \comp \eta_Z = \eta \comp \tilde{\theta}$, by the naturality of the adjunction it must equal $\theta$. 
\end{proof}
\section{Closure Properties of Fibrations}

\begin{definition}\label{cartesian}
	Let $P \colon X \mor Y$ be a functor. A morphism $f \colon x \mor y$ in $X$ is cartesian if given a morphism $h \colon z \mor y$ with $Ph = Pf \comp u$ where $u \colon v \mor Px$ there exists a unique morphism $g \colon z \mor x$ s.t $Pg = u$ and $f \comp g = h$. 
\end{definition}

\begin{definition}
	A functor $P \colon X \mor Y$ is a fibration if for each $x \in X$ and a morphism $f \colon y \mor Px$ in $Y$ there is a cartesian morphism $f' \colon f^{*}x \mor x $ s.t $Pf' \isom f$.
\end{definition}


\begin{theorem}
    Let $P \colon X \to B$ be a fibration and $F \colon Y \mor X $ be an arbitrary functor. Then $F$ itself is a fibration iff 
    \begin{enumerate}
    	\item  $Q \isom P \comp F$ is a fibration and $F$ is a cartesian functor from $Q$ to $P$ and 
    	\item  Each $F_I \colon Y_I \mor X_I$ is a fibration and cartesian morphisms w.r.t these fibrations are stable under reindexing \ie, for every commutative square 
  
    	\[\xymatrix@=0.5cm{
    		Y_1 \ar[rr]^{\varphi_1} \ar[dd]_{\theta} && X_1 \ar[dd]^{\psi} \\
    		\\
    		Y_2 \ar[rr]_{\varphi_2} && Y_1
    		}\]
    	in $Y$ with $\varphi_1$ and $\varphi_2$ cartesian w.r.t $Q$ over $u \colon J \mor I$ and $Q(\theta) = id_J$ and $Q(\psi) = id_I$ it holds that $\theta$ is cartesian w.r.t $F_J$ whenever $\psi$ is cartesian w.r.t $F_I$.
  
    	
    \end{enumerate}
\end{theorem}

\begin{proof}
	($\Longrightarrow$) Since fibration are closed under composition it follows that $Q \is P \comp F$ is a fibration. Consider a $Q-$cartesian morphism $\varphi \colon z \mor y$ over $P(f)$ where $f \colon x \mor Fy$ is $P-$cartesian. Then $PF(\varphi) \isom Q(\varphi) \isom P(f)$, we have two morphisms $F(\varphi)$ and $f$ both to $F(\varphi)$ over $P(f)$, it follows that $F(\varphi) \isom f$. To show $2$, let $y \in Y_I$ and $u \colon x \mor Fy$ in $X_I$. There is a $F-$cartesian morphism $w\colon e \mor e'$ in $Y$ with $F(w) \isom u$. It remains to show that $w \in Y_I$. Since $u \in X_I$ we have $P(w) \isom id_I$ and $PF(w) \isom P(u) \isom id \implies
	w \in Y_I$.
	\\
	($\back$) Let $y \in Y_I$ lying over $x \in X$ ($x \isom F(y) $) and both $x$ and $y$ lie over $b \in B$ ($b \isom P(x) \isom Q(y)$) and morphisms $\delta \colon x' \mor x$ in $X$ and $\beta \colon b' \mor b$ in $B$. We are required to show: there is a cartesian morphism $y' \mor y$ lying over $\delta$. But since $Q$ is a fibration, there is a $Q-$cartesian morphism $\varphi \colon \beta^*y \mor y$ in $Y$ and $F-$cartesian implies that $F(\varphi) \colon \beta^*x \mor x$ is a $P-$cartesian morphism. \\
	Notice that there is a unique factorization of $\delta$ through $F(\varphi)$
	\begin{equation*}
		x' \xlongrightarrow{\alpha} \beta^*x \xlongrightarrow{F(\varphi)} x
	\end{equation*}
	and $\alpha$ is a morphism in $b'$. Since $F_{b'}$ is a fibration, there is a $F_{b'}-$cartesian morphism $\psi \colon \alpha^*\beta^*y \mor \beta^*y$ in $Y_{b'}$. We must show that the composite
	\begin{equation*}
		\varphi \comp \psi \colon \alpha^*\beta^*y \mor \beta^*y \mor y
	\end{equation*}
	is cartesian over $\delta$. Consider the following diagram
	% https://q.uiver.app/#q=WzAsMTMsWzAsMCwiTWFwc19ZKHknLCBcXGFscGhhXipcXGJldGFeKnkpIl0sWzIsMCwiTWFwc19ZKHknLCBcXGJldGFeKnkpIl0sWzQsMCwiTWFwc19ZKHknLCB5KSJdLFswLDIsIk1hcHNfWChGKHknKSwgeCcpIl0sWzIsMiwiTWFwc19YKEYoeScpLCBcXGJldGFeKngpIl0sWzQsMiwiTWFwc19YKEYoeScpLHgpIl0sWzAsNCwiTWFwc19CKFEoeScpLCBiJykiXSxbMiw0LCJNYXBzX0IoUSh5JyksIGInKSJdLFs0LDQsIk1hcHNfQihRKHknKSwgYikiXSxbMywxLCIoMikiXSxbMSwxLCIoMSkiXSxbMSwzLCIoMykiXSxbMywzLCIoMykiXSxbNiw3LCJpZCIsMl0sWzcsOF0sWzQsNV0sWzMsNF0sWzQsN10sWzMsNl0sWzUsOF0sWzIsNV0sWzEsNF0sWzAsM10sWzAsMV0sWzEsMl1d
	\[\xymatrix@=0.5cm{
		Maps_Y(y', \alpha^*\beta^*y) \ar[rr] \ar[dd] && Maps_Y(y', \beta^*y) \ar[rr] \ar[dd] && Maps_Y(y', y) \ar[dd] \\
		& (1) && (2) \\
		Maps_X(F(y'), x') \ar[rr] \ar[dd] && Maps_X(F(y'), \beta^*x) \ar[rr] \ar[dd] && Maps_X(F(y'),x) \ar[dd] \\
		& (3) && (4) \\
		Maps_B(Q(y'), b') \ar[rr]_{\mbox{id}} && Maps_B(Q(y'), b') \ar[rr] && Maps_B(Q(y'), b)
	}\]
	we must show that the composite of the two upper squares is cartesian. To do this, we shall show that both $(1)$ and $(2)$ are cartesian. By our construction $\varphi$ is $Q-$cartesian and $F\varphi)$ is $P-$cartesian, so the composite of $(2)$ and $(4)$ is cartesian and $(4)$ is cartesian, so $(2)$ must be cartesian. To prove that $(1)$ is cartesian, we know that a commutative square is cartesian iff the induced maps on the fibers are equivalences, it suffices to show that the square
	% https://q.uiver.app/#q=WzAsNSxbMCwwLCJNYXBzX1koeScsIFxcYWxwaGFeKlxcYmV0YV4qeSlfayJdLFsyLDAsIk1hcHNfWSh5JywgXFxiZXRhXip5KV9rIl0sWzAsMiwiTWFwc19YKEYoeScpLCB4JylfayJdLFsyLDIsIk1hcHNfWChGKHknKSwgXFxiZXRhXip4KV9rIl0sWzQsMl0sWzIsM10sWzEsM10sWzAsMl0sWzAsMV1d
	\[\xymatrix@=0.5cm{
		Maps_Y(y', \alpha^*\beta^*y)_k \ar[rr] \ar[dd] && Maps_Y(y', \beta^*y)_k \ar[dd] \\
		\\
		Maps_X(F(y'), x')_k \ar[rr] && Maps_X(F(y'), \beta^*x)_k
	}\]
	obtained by taking the fiber at $k \colon Q(y') \mor b'$ is cartesian for all maps $k$. Taking $P-$ and $Q-$cartesian pullbacks along $k$ we can (since $F$ is cartesian functor) identify such a square with the following square
	% https://q.uiver.app/#q=WzAsNSxbMCwwLCJNYXBzX3tZX3tRKHknKX19KHknLCBrXipcXGFscGhhXipcXGJldGFeKnkpIl0sWzIsMCwiTWFwc197WV97USh5Jyl9fSh5Jywga14qXFxiZXRhXip5KSJdLFswLDIsIk1hcHNfe1hfe1EoeScpfX0oRih5JyksIGteKngnKSJdLFsyLDIsIk1hcHNfe1hfe1EoeScpfX0oRih5JyksIGteKlxcYmV0YV4qeCkiXSxbNCwyXSxbMiwzXSxbMSwzXSxbMCwyXSxbMCwxXV0=
	\[\xymatrix@=0.5cm{
		Maps_{Y_{Q(y')}}(y', k^*\alpha^*\beta^*y) \ar[rr] \ar[dd] && Maps_{Y_{Q(y')}}(y', k^*\beta^*y) \ar[dd] \\
		\\
		Maps_{X_{Q(y')}}(F(y'), k^*x') \ar[rr] && Maps_{X_{Q(y')}}(F(y'), k^*\beta^*x)
	}\]
	But this is cartesian since by our assumption the map 
	\begin{equation*}
		k^*\alpha^*\beta^*y \mor k^*\beta^*y
	\end{equation*}
	is $F_{Q(y')}-$cartesian (because  $\psi \colon \alpha^*\beta^*y \mor \beta^*y$ is $F_{b'}-$cartesian).
\end{proof}

\section{Internal Sums}

\begin{definition}\label{Def1}
	Let $\mathcal{B}$ be a category with pullbacks and $P \colon \mathcal{X} \mor \mathcal{B}$ a fibration. A morphism $\varphi \colon X \mor Y$ over $u \colon I \mor J$ is called cocartesian iff for every $v \colon J \mor K$ in $\mathcal{B}$ and $\psi \colon X \mor Z$ over $ v \comp u$ there is a unique morphism $\theta \colon Y \mor Z$ over $v$ with $\psi = \theta \comp \varphi$.
	
	% https://q.uiver.app/#q=WzAsOCxbMiwwLCJYIl0sWzQsMCwiWSJdLFsyLDIsIkkiXSxbNCwyLCJKIl0sWzAsMCwiXFxtYXRoY2Fse1h9Il0sWzAsMiwiXFxtYXRoY2Fse0J9Il0sWzYsMiwiSyJdLFs2LDAsIloiXSxbMCwxLCJcXHZhcnBoaSJdLFsyLDMsInUiLDJdLFs0LDUsIlAiLDJdLFszLDYsInYiLDJdLFswLDcsIlxccHNpIiwwLHsiY3VydmUiOi0zfV0sWzEsNywiXFx0aGV0YSIsMCx7InN0eWxlIjp7ImJvZHkiOnsibmFtZSI6ImRhc2hlZCJ9fX1dXQ==
	\[\xymatrix@=0.5cm{
		X \ar[dd]_{\varphi} \ar[rrdd]^{\psi} &&&& I \ar[dd]_{u} \ar[ddrr] \\
		\\
		Y \ar@{-->}[rr]_{\theta} && Z && J \ar[rr]_{v} && K
	}\] 
	
	The fibration $P$ has internal sums iff the following conditions are satisfied.
	
	\begin{enumerate}
		\item For every $X \in P(I)$ and $u \colon I \mor J$ in $\mathcal{B}$ there is a cocartesian morphism $\varphi \colon X \mor Y$over $u$. \\
		\item The Beck-Chevalley Condition (BCC) holds \ie, for every commuting square in $\mathcal{X}$
		% https://q.uiver.app/#q=WzAsNCxbMCwwLCJDIl0sWzIsMCwiRCJdLFswLDIsIkEiXSxbMiwyLCJCIl0sWzAsMSwiXFx0aWxkZXtcXHZhcnBoaX0iXSxbMSwzLCJcXHBzaSJdLFswLDIsIlxcdGlsZGV7XFxwc2l9IiwyXSxbMiwzLCJcXHZhcnBoaSIsMl1d
		\[\xymatrix@=0.5cm{
			C \ar[rr]^{\tilde{\varphi}} \ar[dd]_{\tilde{\psi}} && D \ar[dd]^{\psi} \\
			\\
			A \ar[rr]_{\varphi} && B
		}\]
		
		over a pullback in $\mathcal{B}$, it holds that $\tilde{\varphi}$ is cocartesian whenever $\varphi$ is cocartesian and $\psi$ and $\tilde{\psi}$ are cartesian.
	\end{enumerate}
\end{definition}

\begin{remark}\label{BCC rem}
	\begin{enumerate}
		\item It is easily seen that for a fibration $\fib$ a morphism $\varphi \colon X \mor Y$ is cocartesian iff for all $\psi \colon X \to Z$ over $P(\varphi)$ there exists a unique vertical morphism $\alpha \colon Y \to Z$ with $\alpha \comp \varphi = \psi$ .
		
		\item One easily sees that the BCC of Definition \ref{Def1} is equivalent to the requirement that for every commuting square in $\mathcal{E}$ 
		% https://q.uiver.app/#q=WzAsNCxbMCwwLCJDIl0sWzIsMCwiRCJdLFswLDIsIkEiXSxbMiwyLCJCIl0sWzIsMywiXFx2YXJwaGkiLDJdLFswLDIsIlxcdGlsZGV7XFxwc2l9IiwyXSxbMSwzLCJcXHBzaSJdLFswLDEsIlxcdGlsZGV7XFx2YXJwaGl9Il1d
		\[\xymatrix@=0.5cm{
			C \ar[rr]^{\tilde{\varphi}} \ar[dd]_{\tilde{\psi}} && D \ar[dd]^{\psi} \\
			\\
			A \ar[rr]_{\varphi} && B
		}\]
		over a pullback in $\mathcal{B}$ it holds that $\psi$ is cartesian whenever $\tilde{\psi}$ is and $\varphi$ and $\tilde{\varphi}$ are cocartesian.
	\end{enumerate}
\end{remark}

We now give an explanation for internal sums in which the Beck-Chevalley Condition arises in a less special way. To this end, we generalize the Fam construction from ordinary categories to fibrations.

\begin{definition}
	Let $\mathcal{B}$ be a category with pullbacks and $\fib$ be a fibration. Then the family fibration $Fam(P)$ for $P$ is defined as $P_{\mathcal{B}} \comp \partial_0^*P$ where
	
	% https://q.uiver.app/#q=WzAsNCxbMCwwLCJQXFxkb3duYXJyb3dcXG1hdGhjYWx7Qn0gIl0sWzIsMCwiXFxtYXRoY2Fse0V9Il0sWzIsMiwiXFxtYXRoY2Fse0J9Il0sWzAsMiwiXFxtYXRoY2Fse0J9XjIiXSxbMCwxXSxbMSwyLCJQIl0sWzMsMiwiXFxwYXJ0aWFsXzAiLDJdLFswLDMsIkZhbShQKSIsMl1d
	\[\xymatrix@=0.5cm{
		{P \downarrow \mathcal{B}} \ar@/_2.5pc/[dddd]_{Fam(P)} \ar[rr] \ar[dd]_{\partial_0^*P} \pullbackcorner && {\mathcal{E}} \ar[dd]^{P} \\
		&
		\\
		{\mathcal{B}^2} \ar[dd]_{P_{\mathcal{B}}} \ar[rr]_{\partial_0} && {\mathcal{B}} \\
		\\
		\mathcal{B}		
	}\]
	is a pullback. The cartesian functor $Fam(P) \colon \partial_0^*P \mor P_{\mathcal{B}}$ is called the fibred family fibration for $P$.
	
	% https://q.uiver.app/#q=WzAsMyxbMCwwLCJQXFxkb3duYXJyb3dcXG1hdGhjYWx7Qn0gIl0sWzIsMCwiXFxtYXRoY2Fse0J9Il0sWzEsMSwiXFxtYXRoY2Fse0J9XjIiXSxbMCwxLCJcXGlvdGFfUCJdLFswLDIsIkZhbShQKSIsMl0sWzIsMSwiUF97XFxtYXRoY2Fse0J9fSIsMl1d
%	\[\xymatrix@=0.5cm{
%		{P \downarrow \mathcal{B}} \ar[rr]^{\iota_P} \ar[dr]_{Fam(P)} && {\mathcal{B}} \\
%		& {\mathcal{B}^2} \ar[ur]_{P_{\mathcal{B}}}
%	}\]
	We can define the cartesian functor $\eta_P \colon P \mor Fam(P)$ as in the diagram 
	
	% https://q.uiver.app/#q=WzAsOSxbMiwyLCJQXFxkb3duYXJyb3dcXG1hdGhjYWx7Qn0gIl0sWzQsMiwiXFxtYXRoY2Fse0V9Il0sWzMsMl0sWzQsNCwiXFxtYXRoY2Fse0J9Il0sWzIsNCwiXFxtYXRoY2Fse0J9XjIiXSxbMiw2LCJcXG1hdGhjYWx7Qn1eMiJdLFswLDQsIlxcbWF0aGNhbHtCfSJdLFswLDAsIlxcbWF0aGNhbHtFfSJdLFsyLDVdLFswLDFdLFsxLDMsIlAiXSxbNCwzLCJcXHBhcnRpYWxfMCIsMl0sWzAsNCwiRmFtKFApIiwyXSxbNCw1LCJQX3tcXG1hdGhjYWx7Qn19Il0sWzYsNCwiXFxEZWx0YV97XFxtYXRoY2Fse0J9fSIsMl0sWzYsNSwiIiwyLHsic3R5bGUiOnsiaGVhZCI6eyJuYW1lIjoibm9uZSJ9fX1dLFs2LDUsIiIsMix7Im9mZnNldCI6LTEsInN0eWxlIjp7ImhlYWQiOnsibmFtZSI6Im5vbmUifX19XSxbNywxLCIiLDIseyJzdHlsZSI6eyJoZWFkIjp7Im5hbWUiOiJub25lIn19fV0sWzcsMSwiIiwyLHsib2Zmc2V0IjotMSwic3R5bGUiOnsiaGVhZCI6eyJuYW1lIjoibm9uZSJ9fX1dLFs3LDYsIlAiLDJdLFs3LDAsIlxcZXRhX1AiLDIseyJzdHlsZSI6eyJib2R5Ijp7Im5hbWUiOiJkYXNoZWQifX19XV0=
	\[\xymatrix@=0.5cm{
		{\mathcal{E}} \ar[dddd]_{P} \ar@{-->}[ddrr]_{\eta_P} \ar@{-}[ddrrrr]  \ar@<0.5ex>@{-}[ddrrrr] \\
		\\
		&& {P\downarrow\mathcal{B}} \ar[rr] \ar[dd]_{\partial_0^*P} \pullbackcorner && {\mathcal{E}} \ar[dd]^{P} \\
		\\
		{\mathcal{B}} \ar@{-}[ddrr] \ar@<0.5ex>@{-}[ddrr] \ar[rr]_{\Delta_{\mathcal{B}}}   && {\mathcal{B}^2} \ar[rr]_{\partial_0} \ar[dd]_{P_{\mathcal{B}}} && {\mathcal{B}} \\
		\\
		&& {\mathcal{B}}
	}\]
	
	where $\Delta_{\mathcal{B}}$ sends a morphism $u \colon I \mor J$ in $\mathcal{E}$ to
	
	% https://q.uiver.app/#q=WzAsNCxbMCwwLCJJIl0sWzIsMCwiSiJdLFswLDIsIkkiXSxbMiwyLCJKIl0sWzAsMSwidSJdLFswLDIsIiIsMCx7Im9mZnNldCI6LTEsInN0eWxlIjp7ImhlYWQiOnsibmFtZSI6Im5vbmUifX19XSxbMiwzLCJ1IiwyXSxbMCwyLCIiLDIseyJzdHlsZSI6eyJoZWFkIjp7Im5hbWUiOiJub25lIn19fV0sWzEsMywiIiwwLHsib2Zmc2V0IjoxLCJzdHlsZSI6eyJoZWFkIjp7Im5hbWUiOiJub25lIn19fV0sWzEsMywiIiwxLHsic3R5bGUiOnsiaGVhZCI6eyJuYW1lIjoibm9uZSJ9fX1dXQ==
	\[\xymatrix@=0.5cm{
		I \ar[rr]^{u} \ar@{-}[dd] \ar@<0.5ex>@{-}[dd] && J \ar@{-}[dd] \ar@<-0.5ex>@{-}[dd] \\
		\\
		I \ar[rr]_{u} && J
	}\]
	in $\mathcal{B}^2$. More explicitly, $\eta_P$ sends $\varphi \colon X \mor Y$ over $u \colon I \mor J$ to a pair
	
	% https://q.uiver.app/#q=WzAsNixbMCwxLCJJIl0sWzIsMSwiSiJdLFswLDMsIkkiXSxbMiwzLCJKIl0sWzAsMCwiWCJdLFsyLDAsIlkiXSxbMCwxLCJ1Il0sWzAsMiwiIiwwLHsib2Zmc2V0IjotMSwic3R5bGUiOnsiaGVhZCI6eyJuYW1lIjoibm9uZSJ9fX1dLFsyLDMsInUiLDJdLFswLDIsIiIsMix7InN0eWxlIjp7ImhlYWQiOnsibmFtZSI6Im5vbmUifX19XSxbMSwzLCIiLDAseyJvZmZzZXQiOjEsInN0eWxlIjp7ImhlYWQiOnsibmFtZSI6Im5vbmUifX19XSxbMSwzLCIiLDEseyJzdHlsZSI6eyJoZWFkIjp7Im5hbWUiOiJub25lIn19fV0sWzQsNSwiXFx2YXJwaGkiXV0=
	\[\xymatrix@=0.5cm{
		X \ar[rr]^{\varphi} && Y \\
		I \ar[rr]^{u} \ar@{-}[dd] \ar@<0.5ex>@{-}[dd] && J \ar@{-}[dd] \ar@<-0.5ex>@{-}[dd] \\
		\\
		I \ar[rr]_{u} && J
	}\]
	in $P\downarrow\mathcal{B}$. Clearly the functor $\eta_P$ preserves the cartesianness of morphisms \ie, $\eta_P$ is cartesian. \vfill
\end{definition}

\begin{remark}
	If a morphism $Fam(P)(\varphi)$ is $P_{\mathcal{B}}$-cocartesian the $\varphi$ is $Fam(P)$-cartesian iff $\varphi$ is $Fam(P)$-cocartesian. Moreover, for every morphism 
	% https://q.uiver.app/#q=WzAsNCxbMCwwLCJBIl0sWzIsMCwiQiJdLFswLDIsIkkiXSxbMiwyLCJKIl0sWzAsMSwidiJdLFswLDIsImEiLDJdLFsyLDMsInUiLDJdLFsxLDMsImIiXV0=
	\[\xymatrix@=0.5cm{
		A \ar[rr]^{v} \ar[dd]_{a} && B \ar[dd]^{b} \\
		\\
		I \ar[rr]_{u} && J
	}\]
	in $\mathcal{B}^2$ we have 

	\[\xymatrix@=0.5cm{
		A \ar@{-}[dd] \ar@<-0.5ex>@{-}[dd] \ar[rr]^v && B \ar@{-}[rr] \ar@<0.5ex>@{-}[rr] \ar@{-}[dd] \ar@<0.5ex>@{-}[dd] && B \ar[dd]^b && A \ar@{-}[dd] \ar@<0.5ex>@{-}[dd] \ar@{-}[rr] \ar@<0.5ex>@{-}[rr]  && A \ar[dd]^a \ar[rr]^v && B \ar[dd]^b \\ &
		1_v && \varphi_b && = && \varphi_a
		\\
		A \ar[rr]_v && B \ar[rr]_b && J && A \ar[rr]_a && I \ar[rr]_u && J
	}\]
	
	where $\varphi_a$ and $\varphi_b$ are $P_{\mathcal{B}}$-cocartesian. From these two observations, one can show that for fibrations $P$ and $Q$ over $\mathcal{B}$ a cartesian functor $F \colon Fam(P) \mor Fam(Q)$ is determined uniquely up to isomorphism by its restriction along the inclusion functor $\Delta_{\mathcal{B}}$ from which it follows that $F \isom Fam({\Delta_B }^*F)$. Thus up to isomorphism all cartesian functors from $Fam(P)$ to $Fam(Q)$ are of the form $Fam(F)$ for some cartesian functor $F \colon P \mor Q$.

\end{remark}

Next we characterize the property of having internal sums in terms of the family monad Fam. \footnote{One can show that $\eta$ and $\mu$ are natural transformations giving rise to a monad (Fam, $\eta$, $\mu$) on \textbf{Fib($\mathcal{B}$)}}.

\begin{theorem}
	Let $\mathcal{B}$ be a category with pullbacks and $\fib$ be a fibration. Then $P$ has internal sums iff $\eta_P \colon P \mor Fam(P)$ has a fibred left adjoint $\coprod_P \colon Fam(P) \mor P$ \ie, $\coprod \dashv \eta_P$ where $\coprod_P$ is cartesian and unit and counit of the adjunction are cartesian natural transformations.
\end{theorem}

\begin{proof}
	($\Longleftarrow$)
	We explicitate the universal property of the unit of the adjunction $\coprod_P \dashv \eta_P$ at an object $(u,X)$ in the following diagram
	% https://q.uiver.app/#q=WzAsOSxbMCwwLCJYIl0sWzIsMCwiWSJdLFs0LDAsIloiXSxbMCwyLCJJIl0sWzIsMiwiSiJdLFs0LDIsIksiXSxbMCw0LCJKIl0sWzIsNCwiSiJdLFs0LDQsIksiXSxbMCwxLCJcXGV0YV97KHUsWCl9Il0sWzMsNCwidSJdLFs0LDUsInYiXSxbMCwyLCJcXHBzaSIsMCx7ImN1cnZlIjotM31dLFsxLDIsIlxcdGhldGEiLDAseyJzdHlsZSI6eyJib2R5Ijp7Im5hbWUiOiJkYXNoZWQifX19XSxbMyw2LCJ1IiwyXSxbNiw3LCIiLDIseyJzdHlsZSI6eyJoZWFkIjp7Im5hbWUiOiJub25lIn19fV0sWzYsNywiIiwyLHsib2Zmc2V0IjotMSwic3R5bGUiOnsiaGVhZCI6eyJuYW1lIjoibm9uZSJ9fX1dLFs0LDcsIiIsMix7InN0eWxlIjp7ImhlYWQiOnsibmFtZSI6Im5vbmUifX19XSxbNCw3LCIiLDIseyJvZmZzZXQiOjEsInN0eWxlIjp7ImhlYWQiOnsibmFtZSI6Im5vbmUifX19XSxbNyw4LCJ2IiwyXSxbNSw4LCIiLDAseyJzdHlsZSI6eyJoZWFkIjp7Im5hbWUiOiJub25lIn19fV0sWzUsOCwiIiwxLHsib2Zmc2V0IjoxLCJzdHlsZSI6eyJoZWFkIjp7Im5hbWUiOiJub25lIn19fV0sWzIsNSwiIiwxLHsic3R5bGUiOnsiYm9keSI6eyJuYW1lIjoiZG90dGVkIn0sImhlYWQiOnsibmFtZSI6Im5vbmUifX19XSxbMSw0LCIiLDEseyJzdHlsZSI6eyJib2R5Ijp7Im5hbWUiOiJkb3R0ZWQifSwiaGVhZCI6eyJuYW1lIjoibm9uZSJ9fX1dLFswLDMsIiIsMSx7InN0eWxlIjp7ImJvZHkiOnsibmFtZSI6ImRvdHRlZCJ9LCJoZWFkIjp7Im5hbWUiOiJub25lIn19fV1d
	\[\xymatrix@=0.5cm{
		X \ar[rr]_{\eta_{(u,X)}} \ar@{.}[dd] \ar@/^1pc/[rrrr]^\psi && Y \ar[rr]_{\theta} \ar@{.}[dd] && Z \ar@{.}[dd]  \\
		\\
		I \ar[rr]^{u} \ar[dd]_{u} && J \ar[rr]^{v} \ar@{-}[dd] \ar@<0.5ex>@{-}[dd]  && K \ar@{-}[dd] \ar@<0.5ex>@{-}[dd]  \\
		\\
		J \ar@{-}[rr] \ar@<0.5ex>@{-}[rr] && J \ar[rr]_{v} && K
	}\]
	
	whose left column is the unit of the adjunction at $(u,X)$. From this it follows that $\eta_{(u,X)} \colon X \mor Y$ is cocartesian over $u$. Cartesianness of $\coprod_P$ says that the cartesian arrow $f$ as given by 
	
	% https://q.uiver.app/#q=WzAsNyxbMCwwLCJYIl0sWzIsMCwiWSJdLFswLDEsIksiXSxbMiwxLCJMIl0sWzAsMywiSSJdLFsyLDMsIkoiXSxbMSwyLCJcXExBUkdFe1xcbHJjb3JuZXJ9Il0sWzAsMSwiXFxwc2kiXSxbMyw1LCJ2Il0sWzQsNSwidSIsMl0sWzIsNCwicCIsMl0sWzIsMywicSJdLFswLDEsIlxcdGV4dHtjYXJ0fSIsMl1d
	\[	\xymatrix@=0.5cm{
		X \ar[rr]_{\mbox{cart}}^\psi && Y \\ 
		K \ar[rr]^{q} \ar[dd]_p  \pullbackcorner &&L \ar[dd]^{v} \\
		&
		\\
		I \ar[rr]_{u} && J
 	}\]
 	

 	
	in $P \downarrow \mathcal{B}$ is sent by $\coprod_P$ to the cartesian arrow $\coprod_Pf$ over $u$ satisfying 
	
	\[ \xymatrix@=0.5cm{
		X \ar[rr]_{\mbox{cart}}^\psi \ar[dd]_{\eta_{(u,X)}} && Y \ar[dd]^{\eta_{(v,Y)}} \\
		\\
		A \ar[rr]_{\coprod_Pf}^{\mbox{cart}} && B
	}
	\]
	where $\eta_{(u,X)}$ and $\eta_{(v,Y)}$  are cocartesian units of the adjunction above $p$ and $v$ respectively. Thus, by $(2)$ of Remark \ref{BCC rem} cartesianness of $\coprod_P$ is just the Beck-Chevalley Condition for internal sums.
	
	$(\Longrightarrow)$ If $P$ has internal sums then the functor $\coprod_P$ left adjoint to $P$ is given by sending a morphism $f$ in $P\downarrow \mathcal{B}$ as given by
	
	\[\xymatrix@=0.5cm{
		X \ar[rr]^{\psi} && Y \\
		K \ar[rr]^{q} \ar[dd]_{p} && L \ar[dd]^{v} \\
		\\
		I \ar[rr]_{u} && J
	}
	\]
	
	to the morphism $\coprod_Pf$ over $u$ satisfying
	\[\xymatrix@=0.5cm{
		X \ar[rr]^{\psi} \ar[dd]_{\varphi_1} && Y \ar[dd]^{\varphi_2} \\
		\\
		A \ar[rr]_{\coprod_Pf} && B
	}\]
	where $\varphi_1$ and $\varphi_2$ are cocartesian over $p$ and $v$ respectively. One can verify that indeed $\coprod_P$ is left adjoint to $\eta_P$ using the units of the adjunction the cocartesian liftings guaranteed for $P$. Again, cartesianness of $\coprod_P$ is easily seen to be equivalent to the Beck-Chevalley Condition.

\end{proof}

The next chapter discusses the dual of internal sums.

\section{Internal Products}
\begin{definition}
    Let $\mathcal{B}$ be a category with pullbacks. A fibration $P \colon \mathcal{X} \mor \mathcal{B}$ has internal products if the following hold:
    \begin{enumerate}
    	\item For $u \colon I \mor J$ in $\mathcal{B}$ and $X \in P(I)$ 
    	
    	\[\xymatrix@=0.5cm{
    		&& u^{*}Z \ar[rr]_{\mbox{cart}}^{\theta} \ar[ddll]_{\alpha}  \ar[dd]^{u^{*}\beta} && Z \ar@{-->}[dd]^{\beta} \\
    		\\
    		X && u^{*}E \ar[ll]^{\varepsilon} \ar[rr]^{\mbox{cart}}_{\varphi}  && E
    	}\]
    	The span $(\varphi, \varepsilon)$ is the determined uniquely up to vertical isomorphism. It is called an evaluation span for $X$ along $u$. \\
    	
    	\item Whenever the pullback
    	
    	\[\xymatrix@=0.5cm{
    		L \ar[rr]^{\tilde{v}} \ar[dd]_{\tilde{u}} \mathlarger{\pullbackcorner} && L \ar[dd]^{u}  \\
    		&
    		\\
    		K \ar[rr]_{v} && J
    	}\]
    	
    	in $\mathcal{B}$ and $(\varphi,\varepsilon)$ is an evaluation span for $X$ along $u$ then 
   
    	with the lower square above the pullback in $\mathcal{B}$, $\tilde{\varepsilon}$ vertical, then ($\tilde{\varphi}, \tilde{\varepsilon}$) is an evaluation span for $\tilde{v}^*X$ along $v$.
    	
    	\[\xymatrix@=0.5cm{
    		\tilde{v}^*Z \ar[rr]_{\mbox{cart}}^{\psi}  && X  \\
    		\\
    		\tilde{u}^*E \ar[rr]_{\mbox{cart}}^{\tilde{\theta}} \ar[uu]^{\tilde{\varepsilon}} \ar[dd]_{\tilde{\varphi}} && u^*E \ar[dd]^{\varphi} \ar[uu]_{u} \\
    		\\
    		\tilde{E} \ar[rr]_{\theta}^{\mbox{cart}} && E
    	}
    	\] 
    	
    \end{enumerate}
\end{definition}

\begin{remark}
	Condition $(2)$ is called the Beck-Chevalley Condition for internal products and essentially says that evaluation spans are stable under reindexing.
\end{remark}

We now give a characterization of internal sums and products in terms of adjunctions from $\S1$ in \cite{JacobsCLTT}.

\begin{lemma}
	Let $\mathcal{B}$ be a category with pullbacks and $\fib$ be a fibration. We say $P$ has internal products (resp. sums) iff the following are satisfied 
	\begin{enumerate}
		\item for every morphism $u \colon I \mor J$ in $\mathcal{B}$, the change of base functor $u^* \colon \mathcal{E}_J \mor \mathcal{E}_I$ has a right adjoint (resp. left adjoint). \\
		
		\item The Beck-Chevalley Condition holds: For every pullback in $\mathcal{B}$ of the form 
		\[\xymatrix@=0.5cm{
			K \ar[rr]^{v} \ar[dd]_{r} \pullbackcorner && L \ar[dd]^{s} \\
			&
			\\
			I \ar[rr]_{u} && J
		}
		\]
		the canonical natural transformation \[
		s^* \prod_u \Longrightarrow \prod_v r^* \qquad \left(\mbox{resp. }  \coprod_{v} r^* \Longrightarrow s^* \coprod_u \right)\] 
		is an isomorphism.
	\end{enumerate}
\end{lemma}

\section{Fibred Terminal Objects}

\begin{definition}
	A fibration $\fib$ is a fibration of categories with terminal object if and only if every fibre $P(I)$ has a terminal object and these are stable reindexing.
\end{definition}

\begin{lemma}\label{terminals}
	Let $\cat{B}$ have a terminal object (besides having pullbacks). A fibration $\fib$ is a fibration of categories with terminal objects if and only if $\cat{E}$ has a terminal $1_{\cat{E}}$ and $P(1_{\cat{E}})$ is terminal in $\cat{B}$ .
\end{lemma}

\begin{proof}
	Suppose $\fib$ is a fibration of categories with terminal objects. Let $1_{\cat{E}}$ be a fibred terminal over $1_\cat{B}$ and $u \colon I \mor 1_\cat{B}$. Let $\bar{u}(1_{\cat{E}}) \colon u^*1_{\cat{E}} \mor 1_{\cat{E}}$ be the cartesian lift of $u$. Then $u^*1_{\cat{E}}$ is a terminal in $P(I)$ and for any $i \in P(I)$ there is a unique vertical arrow $h \colon i \mor u^*1_{\cat{E}}$ with $\bar{u} \comp h$ over $u$. We must show that $1_\cat{E}$ is terminal in $\cat{E}$. Suppose $f \colon i \mor 1_{\cat{E}}$ is another morphism. Then by the universal property of cartesian morphisms, $f = \bar{u} \comp h$. Thus $\bar{u} \comp h$ is the unique arrow from $i$ to $1_{\cat{E}}$. Conversely, suppose $1_{\cat{E}}$ is terminal in $\cat{E}$ with $P(1_\cat{E})$ terminal in $\cat{B}$. Let $u \colon I \mor 1_\cat{B}$ the unique morphism in $\cat{B}$ and $\bar{u}(1_{\cat{E}}) \colon u^*1_{\cat{E}} \mor 1_{\cat{E}}$ be the cartesian lift of $u$. We must show that $u^*1_{\cat{E}}$ is terminal in $P(I)$. To this end, let $a \in P(I)$ and $g \colon a \mor 1_{\cat{E}}$ be the unique morphism in $\cat{E}$. Since $\bar{u}$ is cartesian, there is a unique vertical morphism $l \colon a \mor u^*1_{\cat{E}}$ such that $\bar{u} \comp l  = g$. For any $l^\prime$ with $\bar{u} \comp l^\prime = g$ it follows that $l = l^\prime$ since $1_\cat{E}$ is terminal. This shows that $u^*1_{\cat{E}}$ is terminal in $P(I)$.
\end{proof}
Let $P \colon \cat{E} \mor \cat{B}$ be a fibration. In the category $\slicec{\cat{B}}$, $P$ has a terminal object if $P$ 
\begin{equation*}
	\xymatrix{
		\cat{E} \ar[rr]^{P} \ar[dr]_{P} && \cat{B} \ar[dl]^{id} \\
		& \cat{B}
	}
\end{equation*}
has a right-adjoint. As illustrated below.

\begin{equation*}
	\xymatrix{
		\cat{E} \ar[dr]_{P} && \cat{B} \lltwocell_{P}^{T_I}{'\perp}  \ar[dl]^{id} \\
		& \cat{B}
	}
\end{equation*}

This fibrewise terminal $T_I$ is preserved by cartesianness.

\vspace*{0.2in}
In $\mathbf{Fib}$, $P$ has a terminal object if 

\begin{equation*}
	\xymatrix{
		\cat{E} \ar[d]_{P} \ar[r]^{!} & 1  \ar[d]^{id} \\
		\cat{B} \ar[r]_{!} & 1
	}
\end{equation*}
has a right-adjoint in $\mathbf{Fib}$
\begin{equation*}
	\xymatrix{
		\cat{E} \ar[d]_{P} && 1 \lltwocell_{!}^{T_\cat{E}}{'\perp} \ar[d]^{id}\\
		\cat{B}  && 1 \lltwocell_{!}^{T_\cat{B}}{'\perp} 
	}
\end{equation*}
\ie, both $\cat{E}$ and $\cat{B}$ have terminal objects and $P$ preserves them.

\begin{theorem}
	If $\cat{B}$ has a terminal object, then the following are equivalent
	\begin{enumerate}
		\item $P$ has a terminal object in $\slicec{B}$\\
		\item $P$ has a terminal object in $\mathbf{Fib}$.
	\end{enumerate}
\end{theorem}

\begin{proof}
	By the lifting-and-factorization of adjoints 
	\begin{equation*}
		\xymatrix{
			\cat{E} \ar[drr]_{P} \ar@/^1.5pc/[rrrr]^{!_{\cat{E}}} \ar[rr]^{P} && \cat{B} \ar[d]^{id} \ar[rr]^{!_{\cat{B}}} && 1 \ar[d]^{id} \\
			&& \cat{B} \ar[rr]_{!_{\cat{B}}} && 1
		}
	\end{equation*}
	Given that $!_{\cat{B}} \adj T_\cat{B}$
	\begin{enumerate}[i.]
		\item $P$ has a right-adjoint, $P \adj T_{\cat{E}}$,
		\item $!_{\cat{E}}$ has a right-adjoint $!_{\cat{E}} \adj T_\cat{E}$ over $!_{\cat{B}} \adj T_{\cat{B}}$
	\end{enumerate}
\end{proof}

\section{Fibred Pullbacks}
Let $\cat{B}$ be a category with pullbacks. 
\begin{lemma}
	For a fibration $\fib$ we have
	\begin{enumerate}
		\item A commuting square of cartesian arrows is always a pullback in $\cat{E}$
		\item A commuting square
		\begin{equation*}
			\xymatrix{
				Y_1 \ar[d]_{\beta} \ar[r]^{\varphi_1}_{\text{cart}} & X_1 \ar[d]^{\alpha} \\
				Y_2 \ar[r]_{\varphi_2}^{\text{cart}} & X_2
			}
		\end{equation*}
		in $\cat{E}$ with $\alpha$ and $\beta$ vertical and $\varphi_i$ cartesian is always a pullback.
	\end{enumerate}
\end{lemma}

\begin{proof}(2)
	Consider a cone 
	\[\xymatrix{
		Q \ar[r] \ar[d] & X_1 \ar[d] \\
		Y_2 \ar[r] & X_2
	}\]
	over the cospan $\cospan{Y_2}{X_2}{X_1}{}{}$, then its image is 
	\[\xymatrix{
		P(Q) \ar[dr] \ar[d] \\
		P(Y_2) \ar[r]_{P(\varphi_2)} & P(X_2)
	}\]
	since $P(\varphi_1) = P(\varphi_2)$, another lift is given by 
	\[\xymatrix{
		Q \ar[dr] \ar@{-->}[d] \\
		Y_1 \ar[r]_{\varphi_1} & X_1
	}\]
	then there exists a unique $Q\mor Y_1$ since $\varphi_1$ is cartesian. But we have two \emph{cartesian fillers}, namely, $Q \mor Y_2$ and the $Q \mor Y_1 \mor Y_2$ over 
	\[\xymatrix{
		P(Q) \ar[dr] \ar[d] \\
		P(Y_2) \ar[r]_{P(\varphi_2)} & P(X_2)
	}\]
	since $\varphi_2$ is cartesian, these unique morphisms must be equal. Thus $Q \mor Y_1$ is a cone morphism and unique. Hence the square is a pullback.
\end{proof}
\section{Fibrations of 2-categories}

We recall some basic aspects of fibred 2-categories from \cite{Hermida99}

\begin{definition}
	Let $\fib$ be a 2-functor. A morphism $f \colon X \mor Y$ is said to be cartesian if it is cartesian in a sense of Definition \ref{cartesian} and additionally, for any 2-cells $\sigma \colon h \2cell k \colon Z \mor Y$ and $\gamma \colon v \2cell u \colon PZ \mor PX$, and the lifts $\tilde{v}$ and $\tilde{u}$ of $v$ and $u$ exist, then there is a unique 2-cell $\tilde{\gamma} \colon \tilde{v} \2cell \tilde{u} \colon Z \mor X$.
\end{definition}

\[ \xymatrix{
	Z \ddtwocell^{\tilde{u}}_{\tilde{v}}{\tilde{\gamma}} \ddrrtwocell^{h}_{k}{\alpha}  && && PZ \ddtwocell^{u}_{v}{\gamma} \ddrrtwocell^{Ph}_{Pk}{P\alpha}  \\
	\\ 
	X \ar[rr]_{f} && Y && PX \ar[rr]_{Pf} && PY
}
\]

\section{Fibrant Double Categories}

A double category $\mathbb{D}$ is said to have vertical pullbacks if the structure

\begin{equation*}
	\xymatrix{
		\cat{D}_1 \times_{\cat{D}_0} \cat{D}_1 \ar[r]|-{\odot} & \cat{D}_1 \ar@<1ex>[r]^{s} \ar@<-1ex>[r]_{t} & \cat{D}_0 \ar[l]|-{i}
	}
\end{equation*}
lies in $\mathbf{PBCat}$, the category of categories with pullbacks and pullback-preserving functors. Given a pullback in $\mathcal{D}_0$ we can get the corresponding square of unit of cells in $\mathcal{D}_1$ which is a pullback. $\odot$ preserves pullbacks \ie, we can horizontally compose two vertical pullback squares and obtain a vertical pullback sqaure. This is essentially having a functorial choice of pullbacks.

\begin{definition}[Beck-Chevalley Condition]\label{Double BCC Cell}
	For each pullback of vertical arrows 
	\begin{equation*}
		\xymatrix{
			{} \ar[r]^{q} \ar[d]_{p} \pullbackcorner & {} \ar[d]^{v}\\
			{} \ar[r]_{u} & {}
		}
	\end{equation*}
	the cell 
	\begin{equation*}
		\xymatrix{
			{} \ar[r]|-{\shortmid}^-{p^*} \ar@2{-}[d] \ar@{} [dr]|{\Downarrow} & {} \ar[r]|-{\shortmid}^-{U} \ar[d]_{p} & \ar@{} [dr]|{\Downarrow} {} \ar[d]^{q} \ar[r]|-{\shortmid}^-{q_!} & {} \ar@2{-}[d] \\
			{} \ar[r]|-{\shortmid}_-{U} \ar@2{-}[d] \ar@{} [dr]|{\Downarrow} & {} \ar[d]_{u}  & {} \ar@{} [dr]|{\Downarrow} \ar[r]|-{\shortmid}_-{U}  \ar[d]^{v}& {} \ar@2{-}[d] \\
			{} \ar[r]|-{\shortmid}_-{u_!} & {} \ar[r]|-{\shortmid}_-{U} & {} \ar[r]|-{\shortmid}_-{v^*} & {}
		}
	\end{equation*}
	is invertible.
\end{definition}

\begin{definition}\cite{MR2720187}
	A fibrant double category $\mathbb{D}$ has internal sums whenever the fibration $(s,t) \colon \cat{D}_1 \mor \cat{D}_0 \times \cat{D}_0$ has internal sums \ie, opcartesian lifts of $(s,t)$ satisfy the Beck-Chevalley Condition. 
\end{definition}

\begin{example}\cite{hermida2011categorical}
	$\mathbb{S}$pan($\cat{C}$) as a fibrant double category has internal sums. For a vertical arrow $f \colon X \mor Y$, the extension and restriction along $f$ are given by $\hspan{X}{X}{Y}{id}{f}$ and $\hspan{Y}{X}{X}{f}{id}$
%	\xymatrix@=0.5cm{
%		X & X \ar[l]^{id} \ar[r]_{f} & Y
%	}
%	and 
%	\xymatrix@=0.5cm{
%		Y & X\ar[l]^{f} \ar[r]_{id} & X
%	}
	respectively.
	
	\vspace{0.5cm}
	Note that any span $R \colon \proar{A}{B}$ can be factors as $\delta^* \circ \varepsilon_!$
	where 
	\begin{equation*}
		\delta^* = 
		\vcenter{\vbox{\xymatrix{
					&R \ar[dl]_{\delta} \ar[dr] \\
					A && R
		}}} \qquad
		\varepsilon_! =
		\vcenter{\vbox{\xymatrix{
					&R \ar[dl] \ar[dr]^{\varepsilon} \\
					R && B
		}}}
	\end{equation*}
	More generally we have two embeddings ${}_! \colon \cat{C} \mor \bspan{\cat{C}}$ and ${}^* \colon \cat{C}^{\text{op}} \mor \bspan{\cat{C}}$ such that $f_! \adj f^*$ for $f \colon X \mor Y$. These right adjoints satisfy the Beck-Chevalley Condition 
	\begin{equation*}
		\vcenter{\vbox{\xymatrix{
					{} \ar[r]^{q} \ar[d]_{p} \pullbackcorner & {} \ar[d]^{u} \\
					{} \ar[r]_{v} & {}
		}}} \quad \implies \quad q_! p^* \isom u^* v_!
	\end{equation*}
	\ie
	\begin{equation*}
		\vcenter{\vbox{\xymatrix{
					&{} \ar[dl]_{p} \ar[dr]^{id} && {} \ar[dl]_{id} \ar[dr]^{q} \\
					X && {} && Y 
			}}} \quad
		\isom
		\quad
		\vcenter{\vbox{\xymatrix{
					&{} \ar[dl]_{id} \ar[dr]^{v} && {} \ar[dl]_{u} \ar[dr]^{id} \\
					X && {} && Y 
		}}}
	\end{equation*}
\end{example}

\begin{remark}
	The above span factorization property is inherited by $\mathbb{R}$el as a bicategory. Similarly, we can express a relation $R \colon \proar{X}{Y}$ from $X$ to $Y$ as a span. Composition is given by pullback and 2-cells are just the morphism between the two top objects of the spans. Whence both $\mathbb{R}$el and $\dspan{\cat{C}}$ have internal sums. The Beck-Chevalley Cell in Definition \ref{Double BCC Cell} can be seen from the example above.
\end{remark}


\section{Double Fibred Limits}
Fix $\dcat{B}$ a pseudo category in $\mathbf{PBCat}$, a category of categories with pullbacks and pullback-preserving functors. 

\subsection{Double Terminals}
Consider a double category $\mathbb{D}$ as a pseudo category in $\mathbf{TCat}$, the category of categories with terminal objects. A terminal object $T$ in $\mathbb{D}_0$ gives rise to a terminal object $U_T$ in $\mathbb{D}_1$ such that for any proarrow $\inproar{M}{A}{B}$ there is precisely one cell
\begin{equation*}
	\xymatrix{
		A \ar[d] \ar[r]^{M}|{\shortmid} \ar@{} [dr]|{\theta} & B \ar[d] \\
		T \ar[r]_{U_T}|{\shortmid} & T
	}
\end{equation*}
from $M$ to $U_T$, for any object $M \in \mathbb{D}_1$. This is referred to as the 2-dimensional property of a terminal object. \\ 

\begin{definition}
	A strict double functor $\cat{F} \colon \dcat{E} \mor \dcat{B}$ between pseudo double categories is a double fibration of double categories with terminal objects iff every fibre has a terminal object for $\cat{F}_0$ and $\cat{F}_1$.
\end{definition}

\begin{remark}
	$\cat{F}_0 \colon \dcat{E}_0 \mor \dcat{B}_0$ is an ordinary fibration of categories with terminal objects.
\end{remark}
Let us explicate the fibre at level 1: Let $M \in \dcat{B}_1$ be an object, the fibre $\cat{F}_1(M)$ consists of proarrows $N$ such that $\cat{F}_1N = M$ and cells
\begin{equation*}
	\xymatrix{
		{} \ar[r]^{N}|{\shortmid} \ar@{}[dr]|{\theta} \ar[d] & {} \ar[d] \\
		{} \ar[r]_{L}|{\shortmid} & {}
	}
\end{equation*}
over the globular cell
\begin{equation*}
	\xymatrix{
		{} \ar[r]^{M}|{\shortmid} \ar@{=}[d]  \ar@{}[dr]|{\cat{F}_1\theta} & {} \ar@{=}[d] \\
		{} \ar[r]_{M}|{\shortmid} & {}
	}
\end{equation*}

\begin{remark}
	We shall call cells in the fibre \emph{vertical cells}.
\end{remark}
\begin{definition}
	$\cat{F}_1 \colon \dcat{E}_1 \mor \dcat{B}_1$ is a fibration of categories with terminal objects in the sense that a terminal in the fibre of $M$ is an object $\inproar{U_T}{T}{T}$ such that for any other object $X$ in the fibre is a unique cell 
	
	\begin{equation*}
		\xymatrix{
			{} \ar[r]^{X}|{\shortmid} \ar[d] & {} \ar[d] \\
			T \ar[r]_{U_T}|{\shortmid} & T
		}
	\end{equation*}
	over 
	\begin{equation*}
		\xymatrix{
			{} \ar[r]^{N}|{\shortmid} \ar[d] \ar@{} [dr]|{u} & {} \ar[d] \\
			{} \ar[r]_{M}|{\shortmid} & {}
		}
	\end{equation*}
	
\end{definition}

\begin{proposition}
	Let $\dcat{B}$ also have terminal objects. A strict double fibration $\dfib{F}{E}{B}$ is a double fibration of double categories with terminal objects if and only if $\dcat{E}$ has a terminal $T$ with $FT$ terminal in $\dcat{B}$.
\end{proposition}

\begin{proof}
	Straightfoward application of Lemma \eqref{terminals} in 2-dimensional way.
	
\end{proof}

\subsection{Double Fibred Pullbacks}

\begin{proposition}
	Consider a commuting cube
	\begin{equation*}
		\xymatrix{
			&K \ar@{=}[rr] \ar[dd]|\hole \ar[dl] && K \ar[dd] \ar[dl] \\
			L \ar@{=}[rr] \ar[dd] && L \ar[dd] \\
			& I \ar@{=}[rr]|\hole \ar[dl]  && I \ar[dl]  \\
			J \ar@{=}[rr] && J
		}
	\end{equation*}
	where the top and bottom faces are cartesian cells and the front and back faces are vertical is always a pullback in $\dcat{E}_1$. 	We shall write $(\mathbf{T}), (\mathbf{F}), (\mathbf{B}), (\mathbf{L}), (\mathbf{R}), (\mathbf{G})$ for the top, front, back, left, right, and bottom faces respectively.
\end{proposition}

\begin{proof}
	The $(\mathbf{T})$ and $(\mathbf{G})$ faces are cartesian since the identity functor $i \colon \dcat{E}_0 \mor \dcat{E}_1$ preserves cartesian morphisms. For a double cone $C$ over ($\mathbf{F}$) and ($\mathbf{G}$) we have 
	\begin{equation*}
		\vcenter{\vbox{\xymatrix{
					X \ar[d]_{x_I} \ar[r]|{\shortmid}^{C} \ar@{}[dr]|{\gamma_I} & Y \ar[d]^{y_I} \\					
					I \ar@{=}[r] & I
		}}} \qquad
		\vcenter{\vbox{\xymatrix{
					X \ar[d]_{x_L} \ar[r]|{\shortmid}^{C} \ar@{}[dr]|{\gamma_L} & Y \ar[d]^{y_L} \\					
					L \ar@{=}[r] & L
		}}}
	\end{equation*}
	such that
	\begin{equation*}
		\vcenter{\vbox{\xymatrix{
					X \ar[d]_{x_I} \ar[r]|{\shortmid}^{C} \ar@{}[dr]|{\gamma_I} & Y \ar[d]^{y_I}	 \\
					I \ar[d] \ar@{=}[r] & I \ar[d] \\
					J \ar@{=}[r] & J
		}}} \quad
		=
		\quad
		\vcenter{\vbox{\xymatrix{
					X \ar[d]_{x_L} \ar[r]|{\shortmid}^{C} \ar@{}[dr]|{\gamma_L} & Y \ar[d]^{y_L}	 \\
					L \ar[d] \ar@{=}[r] & L \ar[d] \\
					J \ar@{=}[r] & J
		}}}
	\end{equation*} 
	since ($\mathbf{T}$) is cartesian there is a unique cell 
	
	\begin{equation*}
		\xymatrix@=1cm{
			X \ar@{-->}[d]_{x_K} \ar[r]|{\shortmid}^{C} \ar@{}[dr]|{\tilde{\gamma}_K} & Y \ar@{-->}[d]^{y_K} \\
			K \ar@{=}[r] & K
		}
	\end{equation*}
	there is another unique cell $\tilde{\gamma}_I$ that we started with given by the composition
	
	\begin{equation*}
		\xymatrix@=1cm{
			X \ar[d] \ar[r]|{\shortmid}^{C} \ar@{}[dr]|{\tilde{\gamma}_K} & Y \ar[d] \\
			K \ar[d] \ar@{=}[r] & K \ar[d] \\
			I \ar@{=}[r] & I 
		}
	\end{equation*}
	Since ($\mathbf{G}$)	is cartesian, these unique cells must be equal. Thus the cell $\tilde{\gamma}_K$ is the unique cone morphism.
\end{proof} 

\begin{proposition}
	If $\dfib{P}{E}{B}$ is a double fibration of double categories with pullbacks then every pullback in a fibre is a pullback in $\dcat{E}$. 
\end{proposition}

\begin{proof}
	Clearly, this is true for the fibration $\cat{P}_0$. To show this for $\cat{P}_1$, let
	\begin{equation*}\label{pullbackcube} \tag{$\dagger$}
		\xymatrix@=1cm{
			&L \ar@{=}[rr] \ar[dd]|\hole \ar[dl]_{\tilde{v}} && L \ar[dd]^{k} \ar[dl]_{\tilde{v}} \\
			I \ar@{=}[rr] \ar[dd]_{u} && I \ar[dd] \\
			& K \ar@{=}[rr]|\hole \ar[dl]_{v}  && K \ar[dl]_{v}  \\
			J \ar@{=}[rr] && J
		}
	\end{equation*}
	be a pullback in $\dcat{E}_1$ over  $\inproar{M}{}{} \in \dcat{B}_1$. Consider the following diagram 
	\begin{equation*}
		\xymatrix@=1cm{
			X \ar@/_3pc/[ddd]_{x_K} \ar[d]_{x_I} \ar[r]|{\shortmid}^{C} & Y \ar[d]^{y_I} \ar@/^3pc/[ddd]^{y_K} \\
			I \ar[d]_{u} \ar@{=}[r] & I \ar[d]^{u} \\
			J \ar@{=}[r] & J \\
			K \ar[u]^{v} \ar@{=}[r] & K \ar[u]_{v}
		}
	\end{equation*}
	with $v \comp x_K = u \comp x_I$ and $v \comp y_K = u \comp y_I$ for an object ${C} \in \dcat{E}_1$ \ie, a cone over ($\mathbf{F}$) and ($\mathbf{G}$). The cells 
	\begin{equation*}
		\vcenter{\vbox{\xymatrix{
					X \ar[d]_{x_K} \ar[r]|{\shortmid}^{C} & Y \ar[d]^{y_K} \\
					K \ar@{=}[r] & K
		}}}
		\qquad
		\vcenter{\vbox{\xymatrix{
					X \ar[d]_{x_I} \ar[r]|{\shortmid}^{C} & Y \ar[d]^{y_I} \\
					I \ar@{=}[r] & I
		}}}
	\end{equation*}
	lie over the same cell $\theta \colon N \mor M$ in $\dcat{B}_1$. Let the cells 
	\begin{equation*}
		\vcenter{\vbox{\xymatrix{
					f^*K \ar[d]_{\cleav{f}{K}} \ar@{=}[r] & f^*K \ar[d]^{\cleav{f}{K}} \\
					K \ar@{=}[r] & K
		}}}
		\qquad
		\vcenter{\vbox{\xymatrix{
					f^*I \ar[d]_{\cleav{f}{I}} \ar@{=}[r] & f^*I \ar[d]^{\cleav{f}{I}} \\
					I \ar@{=}[r] & I
		}}}
	\end{equation*}
	be cartesian \ie, the vertical arrows are cartesian lifts of $f$ in the base. Moreover, let 
	\begin{equation*}
		\vcenter{\vbox{\xymatrix{
					X \ar[d] \ar[r]|{\shortmid}^{C} & Y \ar[d] \\
					f^*K \ar@{=}[r] & f^*K
		}}}
		\qquad
		\vcenter{\vbox{\xymatrix{
					X \ar[d] \ar[r]|{\shortmid}^{C} & Y \ar[d] \\
					f^*I \ar@{=}[r] & f^*I
		}}}
	\end{equation*}
	be vertical cells over $\inproar{N}{}{}$. We know that \eqref{pullbackcube} is a pullback under change of base, there is a vertical cell
	\begin{equation*}
		\xymatrix@=1cm{
			X \ar[d] \ar[r]|{\shortmid}^{C} & Y \ar[d] \\
			f^*L \ar@{=}[r] & f^* L
		}
	\end{equation*}
	with 
	\begin{equation*}
		\vcenter{\vbox{\xymatrix{
					X \ar[d] \ar[r]|{\shortmid}^{C} & Y \ar[d] \\
					f^*K \ar@{=}[r] & f^*K
		}}}
		\quad
		=
		\quad
		\vcenter{\vbox{\xymatrix{
					X \ar[d] \ar[r]|{\shortmid}^{C} & Y \ar[d] \\
					f^*L \ar[d] \ar@{=}[r] & f^*L \ar[d] \\
					f^*K \ar@{=}[r] & f^*K
		}}}
	\end{equation*}
	similarly for $f^*I$. We obtain the following diagram 
	\begin{equation*}
		\xymatrix@=1cm{
			X  \ar[d] \ar[rrrrrr]|{\shortmid}^{C}  &&&&&& Y \ar[d] \\
			f^*L \ar@{=}[rr] \ar[dr] \ar[dd] && f^*L \ar[dr] \ar[dd]|\hole \ar@{=}[rr] && f^*L \ar[dr] \ar[dd]|\hole \ar@{=}[rr] && f^*L \ar[dr]\ar[dd]|\hole \\
			& L \ar@{=}[rr] \ar[dd] && L \ar[dd] \ar@{=}[rr] && L \ar[dd] \ar@{=}[rr] && L \ar[dd] \\
			f^*K \ar@{=}[rr]|\hole \ar[dr] \ar[dd] && f^*K \ar[dd]|\hole  \ar[dr] && f^*I \ar[dr] \ar[dd] \ar@{=}[rr]|\hole && f^*I \ar[dr] \ar[dd]|\hole \\
			& K \ar@{=}[rr] \ar[dd] && K \ar[dd]  && I \ar[dd] \ar@{=}[rr] && I \ar[dd] \\
			f^*J \ar@{=}[rr]|\hole \ar[dr] && f^*J \ar[dr] \ar@{=}[rr]|\hole && f^*J \ar[dr] \ar@{=}[rr]|\hole && f^*J \ar[dr] \\
			& J \ar@{=}[rr] && J \ar@{=}[rr] && J \ar@{=}[rr] && J
		}
	\end{equation*}
	where the cells 
	\begin{equation*}
		\vcenter{\vbox{\xymatrix{
					f^*L \ar[d]_{\cleav{f}{L}} \ar@{=}[r] \ar@{}[dr]|{\text{cart}} & f^*L \ar[d]^{\cleav{f}{L}} \\
					L \ar@{=}[r] & L
		}}}
		\qquad
		\vcenter{\vbox{\xymatrix{
					f^*K \ar[d]_{\cleav{f}{K}} \ar@{=}[r] \ar@{}[dr]|{\text{cart}} & f^*K \ar[d]^{\cleav{f}{K}} \\
					K \ar@{=}[r] & K
		}}}
		\qquad
		\vcenter{\vbox{\xymatrix{
					f^*I \ar[d]_{\cleav{f}{I}} \ar@{=}[r] \ar@{}[dr]|{\text{cart}} & f^*I \ar[d]^{\cleav{f}{I}} \\
					I \ar@{=}[r] & I
		}}}
	\end{equation*}
	are cartesian. The desired cell is given by 
	\begin{equation*}
		\vcenter{\vbox{\xymatrix{
					X \ar[d] \ar[r]|{\shortmid}^{C} \ar@{}[dr]|{\xi} & Y \ar[d] \\
					L \ar@{=}[r] & L
		}}}
		\qquad
		=
		\qquad
		\vcenter{\vbox{\xymatrix{
					X \ar[d] \ar[r]|{\shortmid}^{C}  \ar@{}[dr]|{\gamma}& Y \ar[d] \\
					f^*L \ar[d]_{\cleav{f}{L}} \ar@{=}[r] \ar@{}[dr]|{\psi} & f^*L \ar[d]^{\cleav{f}{L}} \\
					L \ar@{=}[r] & L
		}}}
	\end{equation*}
	For any other cell $\xi' = \frac{\gamma'}{\psi}$. It follows that $\gamma' = \gamma$ as both are mediating arrows to the pullback under change of base for the cone given by vertical cells into $f^*K$ and $f^*I$ and therefore, it follows that $\xi = \xi'$.
	
\end{proof}

\subsection{Tabulators and Fibred Tabulators}
\begin{definition}\label{tabulator}
	In a double category $\dcat{D}$, a \emph{1-tabulator} of a proarrow $M$ is an object $1^{\bullet}_{T}$ together with a cell
	\begin{equation*}
		\xymatrix@=1cm{
			T \ar@{=}[r] \ar[d] \ar@{}[dr]|{\pi} & T \ar[d] \\
			{} \ar[r]|{\shortmid}_M & {}
		}
	\end{equation*}
	such that for any other cell 
	\begin{equation*}
		\xymatrix{
			R \ar@{=}[r] \ar[d] \ar@{}[dr]|{\pi \prime} & R \ar[d] \\
			{} \ar[r]|{\shortmid}_M & {}
		}
	\end{equation*}
	we have 
	\begin{equation*}
		\vcenter{\vbox{\xymatrix{
					R \ar@{=}[r] \ar[d] \ar@{}[dr]|{\pi \prime} & R \ar[d] \\
					{} \ar[r]|{\shortmid}_M & {}
		}}}
		\quad
		=
		\quad
		\vcenter{\vbox{\xymatrix{
					R \ar@{=}[r] \ar@{-->}[d]_k \ar@{}[dr]|{U_{k}} & R \ar@{-->}[d]^k \\
					T \ar@{=}[r] \ar[d] \ar@{}[dr]|{\pi} & T \ar[d] \\
					{} \ar[r]|{\shortmid}_M & {}
		}}}
	\end{equation*}
\end{definition}


\begin{lemma}[Grandis-Parè]
	A double category $\dcat{D}$ has 1-tabulators if and only if the functor $i \colon \dcat{D}_0 \mor \dcat{D}_1$ has a right adjoint. 
\end{lemma}

We shall refer to 1-tabulators with tabulators only.

\begin{definition}
	Let $\dcat{B}$ be a double category with tabulators. A strict double fibration $\dfib{P}{E}{B}$ is double fibration of double categories with tabulators if and only if every fibre has tabulators and they are stable reindexing.
\end{definition}

\begin{lemma}\label{carttabulator}
	For a double fibration $\dfib{P}{E}{B}$, a cartesian cells of the form 
	\begin{equation*}
		\xymatrix@=1cm{
			{}  \ar[d] \ar@{=}[r] \ar@{}[dr]|{\text{cart}} & {} \ar[d] \\
			{} \ar@{=}[r] & {}
		}
	\end{equation*}
	can be reduced to tabulators of proarrow identities in $\dcat{E}$. 
\end{lemma}
\begin{proof}
	Immediate from Definition \ref{tabulator}.
\end{proof}

\begin{lemma}\label{localtabulator}
	For a double fibration $\dfib{P}{E}{B}$ of double categories with tabulators, a tabulator in a fibre is a tabulator in $\dcat{E}$.
\end{lemma}

\begin{proof}
	Let 
	\begin{equation*}
		\xymatrix@=1cm{
			Tu \ar@{=}[r] \ar[d]_{l} \ar@{}[dr]|{\pi_u} & Tu \ar[d]^{k} \\
			A \ar[r]|{\shortmid}_u & B
		}
	\end{equation*}
	be a tabulator over $1^{\bullet}_{M}$  and $1^{\bullet}_{f} \colon 1^{\bullet}_{N} \mor 1^{\bullet}_{M}$ be a cell in $\dcat{B}$ . Consider a double cone over $\pi_u$ given by 
	\begin{equation*}
		\xymatrix@=1cm{
			&C \ar[dl]_{c_T} \ar[ddl]|\hole^(.59){h_1} \ar[dr]^{c_T} \ar[ddr]|\hole_(.59){h_2}\\
			Tu \ar[d]_{l} \ar@{=}[rr] && Tu \ar[d]^{k} \\
			A \ar[rr]|{\shortmid}_u && B
		}
	\end{equation*}
	with $k \comp c_T = h_2$ and $l \comp c_T = h_1$. Clearly, the cell 
	\begin{equation*}
		\xymatrix@=1cm{
			C \ar@{=}[r] \ar[d]_{h_1} \ar@{}[dr]|{\theta} & C \ar[d]^{h_2} \\
			A \ar[r]|{\shortmid}_u & B
		}
	\end{equation*}
	lies over $1^{\bullet}_{f}$ in $\dcat{B}$.
	Let 
	\begin{equation*}
		\xymatrix@=1cm{
			C \ar[d]_{h^*_1} \ar@{=}[r] \ar@{}[dr]|{\theta^*} & C \ar[d]^{h^*_2} \\
			f^*A \ar[d]_{\cleav{f}{A}} \ar[r]|{\shortmid}_{u^*} \ar@{}[dr]|{\text{cart}} & f^*B \ar[d]^{\cleav{f}{B}} \\
			A \ar[r]|{\shortmid}_{u} & B
		}
	\end{equation*}
	be the vertical-cartesian factorization of $\theta$. Since every fibre has a tabulator, there is a vertical cell
	\begin{equation*}
		\xymatrix@=1cm{
			C \ar[d]_{c^*_T} \ar@{=}[r] \ar@{}[dr]|{\gamma} & C \ar[d]^{c^*_T} \\
			f^*Tu \ar@{=}[r] & f^*Tu
		}
	\end{equation*}
	such that
	\begin{equation*}
		\vcenter{\vbox{\xymatrix{
			C \ar[d]_{c^*_T} \ar@{=}[r] \ar@{}[dr]|{\gamma} & C \ar[d]^{c^*_T} \\
			f^*Tu \ar[d]_{l^*} \ar@{=}[r] \ar@{}[dr]|{\pi^*_u} & f^*Tu \ar[d]^{h^*}	\\	
			f^*A \ar[r]|{\shortmid}_{u^*} & f^* B
		}}}
		\quad
		=
		\quad	
		\vcenter{\vbox{\xymatrix{
			C \ar[d]_{h^*_1} \ar@{=}[r] \ar@{}[dr]|{\theta^*} & C \ar[d]^{h^*_2} \\
			f^*A \ar[r]|{\shortmid}_{u^*} & f^* B	
		}}}
	\end{equation*}
	We obtain the following diagram
	\begin{equation*}
		\xymatrix@=1cm{
			&& C \ar[ddrr]^{h^*_2} \ar[ddll]_{h^*_1} \ar[d]_{\gamma} \\
			&& f^*Tu \ar[dll]^{l^*} \ar[drr]_{k^*} \ar[dd]|\hole_(.20){\xi}
			\\
			f^*A \ar[dd]_{\cleav{f}{A}}\ar[rrrr] &&&& f^*B \ar[dd]^{\cleav{f}{B}}\\
			&& Tu \ar[dll]_{l} \ar[drr]^{k} \\
			A \ar[rrrr]|{\shortmid}_{u} &&&& B
		}
	\end{equation*}
	The cell 
	\begin{equation*}
		\vcenter{\vbox{\xymatrix{
					C \ar[d]_{c_T} \ar@{=}[r] \ar@{}[dr]|{\omega} & C \ar[d]^{c_T} \\
					Tu \ar@{=}[r] & Tu
		}}}
		\quad
		=
		\quad
		\vcenter{\vbox{\xymatrix{
					C \ar[d]_{c^*_T} \ar@{=}[r] \ar@{}[dr]|{\gamma} & C \ar[d]^{c^*_T} \\
					f^*Tu \ar@{=}[r] \ar[d]_{t} \ar@{}[dr]|{\psi} & f^*Tu \ar[d]^{t} \\
					Tu \ar@{=}[r] & Tu
		}}}
	\end{equation*}
	is the desired mediating cell.
\end{proof}

I have made this interesting observation:
\begin{lemma}\label{pulltabulators}
	For a a strict double fibration $\dfib{P}{E}{B}$ of double categories with tabulators. For a commuting cube
	\begin{equation*}
		\xymatrix@=1cm{
			&{} \ar[rr] \ar[dd]|\hole \ar[dl]_{\tilde{v}} && {} \ar[dd]^{k} \ar[dl]_{\tilde{v}} \\
			{} \ar[rr] \ar[dd]_{u} && {} \ar[dd] \\
			& {} \ar[rr]|\hole \ar[dl]_{v}  && {} \ar[dl]_{v}  \\
			{} \ar[rr] && {}
		}
	\end{equation*}
	given by $(\mathbf{T}), (\mathbf{F}), (\mathbf{B}), (\mathbf{L}), (\mathbf{R}), (\mathbf{G})$ as the top, front, back, left, right, and bottom faces respectively, we have the following
	\begin{enumerate}
		\item If ($\mathbf{B}$) is tabulator cell and ($\mathbf{G}$) is cartesian, then ($\mathbf{T}$) is cartesian
		\item If ($\mathbf{F}$) is a tabulator and ($\mathbf{G}$) is cartesian then ($\mathbf{B}$) is a tabulator if and only if ($\mathbf{T}$) is cartesian.
		\item If ($\mathbf{T}$) and ($\mathbf{B}$) are tabulator cells and ($\mathbf{G}$) is cartesian, then the cube is a pullback.
	\end{enumerate}
\end{lemma}

\begin{remark}
	Condition 2 in Lemma \ref{pulltabulators} is a Beck-Chevalley like condition for fibred tabulators. 
\end{remark}

\begin{remark}
	In the above formulation, we require the source $s$ and target $t$ maps to preserve the fibre structure. Moreover, each property of a double fibration (\ie, terminal objects or pullbacks) should lie in the 2-category of fibrations with terminal objects or pullbacks. Since tabulators only appear in the double categorical case, the treatment is purely double categorical. 
\end{remark}

Since we gave a characterization theorem for double terminals, we also have the following characterization theorems for pullbacks and tabulators.

\begin{theorem}
	A double fibration $\dfib{P}{E}{B}$ is a double fibration of double categories with pullbacks if and only if $\dcat{E}$ has and $\cat{P}$ preserve them.
\end{theorem}
This follows readily from the double fibration structure since it holds for $\cat{P}_0$. However, the following theorem does not.

\begin{theorem}
	Let $\dcat{B}$ have tabulators. A double fibration $\dfib{P}{E}{B}$ is a double fibration of double categories with tabulators if and only if $\dcat{E}$ has tabulators  and $\cat{P}$ preserves them. 
\end{theorem}
\begin{proof}
	Let 
	\begin{equation*}
		\xymatrix@=1cm{
			A \ar@{=}[r] \ar[d]_{g} & A \ar[d]^{g}\\
			B \ar@{=}[r] & B
		}
	\end{equation*}
	be a cell in $\dcat{E}$. Consider the following diagram
	\begin{equation*}
		\xymatrix@=1cm{
			A \ar@{=}[r] \ar[d]_{a} \ar@{}[dr]|{\pi_{f^*B}} & A \ar[d]^{a}\\
			f^*B \ar@{=}[r] \ar[d]_{\cleav{f}{B}} \ar@{}[dr]|{\epsilon} & f^*B \ar[d]^{\cleav{f}{B}} \\
			B \ar@{=}[r] & B
		}
	\end{equation*}
	where $\epsilon$ is cartesian over $1^{\bullet}_{f} \colon 1^{\bullet}_{N} \mor 1^{\bullet}_{M}$ and $\pi_{f^*B}$ is a tabulator in the fibre. Then $\pi_{f^*B}$ is a tabulator in $\dcat{E}$ by Lemma \ref{localtabulator} and $\epsilon$ is a tabulator in $\dcat{E}$ by Lemma \ref{carttabulator}. It remains to see that the composite $\frac{\pi_{f^*B}}{\epsilon}$ is a tabulator in $\dcat{E}$. 
	
	
\end{proof}



\section{Dual Notions}
\subsection{2-Opfibrations}
\begin{definition}
	A 2-functor $P \colon \cat{E} \mor \cat{B}$ is a 2-opfibration if the following hold
	\begin{enumerate}
		\item For any $e \in \cat{E}$  and $u \colon Pe \mor b$, there is an opcartesian lift $\bar{u} \colon e \mor u_!e$ with $P\bar{u} = u$ ;
		\item for any $g \in \cat{E}$ and $\alpha \colon Pg \2cell h$, there is an opcartesian lift $\bar{\alpha} \colon g \2cell \alpha_!g$ with $P\bar{\alpha} = \alpha$; and 
		\item Horizontal composition of opcartesian 2-cells is again opcartesian.
	\end{enumerate}
\end{definition}

\begin{remark}
	The above definition refers to coop-2-fibrations as 2-opfibrations. Comparing to Hermida, opcartesian 1-cells and 2-cells would be defined using pre-composition instead of post-composition.
\end{remark}

\subsection{Double Bifibrations}
A strict double functor $\mathscr{P} \colon \mathbb{E} \mor \mathbb{B}$ defines a double fibration if the components $P_0$ and $P_1$ are fibrations. We say $\mathscr{P}$ is a double fibration if in addition the components $P_0$ and $P_1$ are opfibrations. 
\subsubsection{Relation to 2-fibrations}
\begin{proposition}
	A 2-functor $\mathscr{P} \colon \cat{E} \mor \cat{B}$ is a 2-fibration if and only if $Q(\mathscr{P}) \colon Q(\cat{E}) \mor Q(\cat{B})$ is a double fibration. 
\end{proposition}

\begin{proof}($\2cell$) 
	The cartesian 1-cells make the functor $Q(P_0)=|P| \colon |\cat{E}| \mor |\cat{B}|$ an ordinary fibration. We are required to show: $Q(P)_1$ is a fibration and the external source and target maps of $Q(\cat{E})$ are cleavage-preserving. To this end suppose there is an object $\proar{X}{Y}$ in $Q(\cat{E})_1$ and a morphism
	\begin{equation*}
		\xymatrix{
			A \ar[r]^{u} \ar[d]_{f} \ar@{}[dr]|{\Downarrow \alpha} & B \ar[d]^{g} \\
			PX \ar[r]_{Pv} & PY
		}
	\end{equation*}
	in $Q(\cat{B})_1$. The lift of $\alpha$ in $Q(\cat{E})_1$ can be constructed as follows: Since $\mathscr{P}$ is locally a fibration, the cartesian lift $\tilde{\alpha}$ of $\alpha$ exists. Since $\sigma(g,Y)$ is cartesian, there is a unique $\tilde{u} \colon f^*X \mor g^*Y$ such that the following triangle commutes
	
	\begin{equation}\label{cart2cell} \tag{$\clubsuit$}
		\xymatrix{
			f^*X \ar@{-->}[r]^{\tilde{u}} \ar[d]_{\sigma(f,X)} \ar@{}[dr]|{\Downarrow \tilde{\alpha}} & g^*Y \ar[d]^{\sigma(g,Y)}\\
			X \ar[r]_v & Y
		}
	\end{equation}
	Claim: \eqref{cart2cell} is cartesian w.r.t $Q(P)_1$
	\begin{proof}
		Since the vertical arrows are cartesian and the 2-cell is cartesian, this means $\eqref{cart2cell}$ will be $Q(P)_1-$cartesian.
	\end{proof}

	
	
\end{proof}

\section{Double Internal Products}

Double Evaluation Spans

\begin{equation*}
	\xymatrix{
		X \ar@{=}[rr] && X \\
		& u^*Z \ar@{=}[rr] \ar[dd]|\hole \ar[dl] \ar[lu] && u^*Z \ar[lu] \ar[dl] \ar[dd]  \\
		u^*E \ar@{=}[rr] \ar[dd] \ar[uu] && u^*E \ar[dd]|\hole \ar[uu]|\hole \\
		& Z  \ar@{=}[rr] \ar@{-->}[dl] && Z  \ar@{-->}[dl]\\
		E \ar@{=}[rr] && E
	}
\end{equation*}





%% Maybe I need to add one more part: Examples.
%% Set style and colour later.
%
%\subsection{Pictures}
%
%\begin{figure}[htbp]
%    \center
%    \includegraphics[scale=0.06]{img/photo.jpg}
%    \caption{Sydney, NSW}
%\end{figure}

%\subsection{Citation}
%
%This is a citation\cite{Eg}.

\newpage

% ------------------------------------------------------------------------------
% Reference and Cited Works
% ------------------------------------------------------------------------------

\bibliographystyle{IEEEtran}
\bibliography{references.bib}


% ------------------------------------------------------------------------------

\end{document}
